\section{The marked higher dictionary on the prescribed coefficient field}
\label{sec:prescribed-higher-dictionary}

The higher marking changes the integral period lattices. To read the
character on the resulting moduli scheme one needs a period chart
for each adelic lattice, using the given embedding of the lower
lattice and the given archimedean frame. We first supply the two
section laws used to normalize those charts.

% BEGIN D74U1
\paragraph{D74U1. Sections of a polarized geometric fiber.}
\label{claim:D74U1}
Let $P$ be an original $\mathrm{PolAV}$ over a scheme $S$. On every
geometric fiber $B$ of dimension $g=P.g$, write $L$ for the pulled
line. For every $q\geq1$ and $c\in B(k)$, both $L^q$ and
$N=L^{q-1}\otimes t_c^*L$ have $q^g$ independent global sections
forming a basis over the algebraically closed residue field $k$, and
their higher cohomology vanishes.

\begin{proof}
PBC25 transports the polarization to the fiber, including the field
$\mathrm{KTriv}$. In particular $L$ is ample and symmetric.
Consider the polarization isogeny
$\lambda_L:B\to B^\vee$, $a\mapsto[t_a^*L\otimes L^{-1}]$.
We check its scheme-theoretic kernel, which matters in positive
characteristic. For every affine $T/k$ and $a\in\ker\lambda_L(T)$,
put $D_a=t_a^*L_T\otimes L_T^{-1}$. Rigidify $D_a$ by tensoring
with the inverse of the pullback of its identity restriction.
The rigidified bundle then has zero relative Picard class. It has a
unique rigidified trivialization fppf locally on $T$: membership
in the zero class gives a pullback from the base locally, and
rigidifying such a pullback trivializes it. Uniqueness follows
from $(B_T\to T)_*\mathcal O_{B_T}=\mathcal O_T$. For this last
identity over arbitrary $T$, one may work over an affine of $T$
and use flat base change from the field $k$ and
$H^0(B,\mathcal O_B)=k$. Uniqueness descends the trivialization.
Thus $D_a$ is the pullback of a line on $T$ and has the generators
Zariski locally on $T$ required by $\mathrm{KTriv}$. The latter
gives $a=0$. Affine tests include the possibly nonreduced kernel
itself. The isogeny has degree one.

The degree formula and Riemann--Roch give
$\deg\lambda_L=\chi(L)^2$ and $\chi(L^q)=q^g\chi(L)$.
Ampleness gives positive top self-intersection in positive dimension,
so $\chi(L)=1$. In dimension zero $B=\operatorname{Spec} k$
and the formula holds directly. For vanishing, choose arbitrarily
large integers $b$ invertible in $k$. The pullback
$[b]^*L^q\simeq L^{qb^2}$ follows from symmetry and the theorem of
the cube. The unit followed by the trace for $[b]$ multiplies by
the invertible degree $b^{2g}$. After applying the projection formula,
$H^j(B,L^q)$ is a summand of $H^j(B,L^{qb^2})$. Serre vanishing
therefore gives the assertion for $j>0$.

Choose $a\in B(k)$ with $qa=c$, using surjectivity of multiplication
on the abelian variety, also when the characteristic divides $q$.
The theorem of the square gives
$\lambda_{L^q}(a)=q\lambda_L(a)=\lambda_L(c)$, hence
$N\simeq t_a^*L^q$. Translation preserves cohomology. This choice
is needed only on an algebraically closed fiber.
The classical degree and Riemann--Roch formulas are
J.~S.~Milne, \emph{Abelian Varieties}, Theorem~13.3(a,b),
\url{https://www.jmilne.org/math/xnotes/AVs.pdf}.
\end{proof}
% END D74U1

% BEGIN D74U2
\paragraph{D74U2. The section bases before every affine test.}
\label{claim:D74U2}
For the actual polarized data the geometric rank,
universal section bases and residue lifts satisfy all four native laws:
\[
 \begin{gathered}
 \mathrm{GlobalSectionsRankAlgClosed},\qquad \mathrm{SecBaseChangeAnySt},\\
 \mathrm{SecBaseChangeSt},\qquad \mathrm{SecTwistSt}.
 \end{gathered}
\]

\begin{proof}
The absolute rank is D74U1 with the identity pullback and scalar
identifications of SGV2. For the bases let $p:B\to S'$ be the
original abelian scheme, or its pullback to the affine local test
of the twist law. Write $N$ for the indicated power or translated
power there. The map $p$ is proper, flat and of finite presentation
by the abelian-scheme fields. The line $N$ is of finite presentation
and flat over $S'$. The proper-support perfect cohomology theorem
makes $Rp_*N$ perfect with formation commuting with arbitrary derived
base change (Stacks Project, Lemma 36.30.7,
\url{https://stacks.math.columbia.edu/tag/0A1J}).
In particular its derived geometric fibers are the complexes computed
in D74U1.

At any point choose a finite free complex representing $Rp_*N$
locally. Split each unit pivot of each differential at the point as
a two-term identity summand, shrinking to invert its pivot. In the
remaining complex all differentials are zero modulo the maximal ideal.
After passage to the algebraic closure of the residue field, D74U1
therefore identifies the rank in each degree with the dimension of
the cohomology in that degree. Only degree zero remains and its rank
is $q^g$. The construction takes finitely many pivots, so gives one
open neighborhood and a basis there before choosing any later
affine test. In degree zero the derived comparison is the usual
pullback of sections, by functoriality. Thus on any affine test the
pulled basis spans and is independent, including over a nonflat
algebra. By the chart gluing in GSL1, every list satisfying the
original $\mathrm{IsPullSec}$ consists of these same pullbacks.
This proves the arbitrary-scheme law and its locally Noetherian
restriction, for every positive power.

For the translated line on $\operatorname{Spec}R$ with $R$ local,
the resulting finite projective module is free of rank $q^g$. Its
basis gives unique coordinates for each section on each of the
source's common charts, precisely $\mathrm{IsSecBasis}$. An arbitrary
section on the residue field has a coordinate vector there by base
change. Lift each of its $q^g$ coordinates to $R$ and form that
linear combination over $R$. Its pullback has the required
coordinates on every common chart by GSL1, giving the specified
$\mathrm{IsPullSecLC}$ equality. This proves both clauses of the
twist law.
\end{proof}
% END D74U2

% BEGIN D74U3
\paragraph{D74U3. Upper period models with the specified lower frame.}
\label{claim:D74U3}
Fix a slot $i$, put $V=W_c(D,i)$, $r=D.R.r$ and $D_w=T.Q.Dw(i+1)$.
Take the prescribed level, upper moduli object and coefficient record
of A73H2 and L73D2. Write $M=c'.F=c_2.F$ as an intermediate field
of $\mathbb C$. Use the original $\iota=\mathrm{iotaW}(T,i)$ and
$x_0=T.x_0(i+1,\mathrm{sOv}(T,c',\mathrm{hle}))$. There is a finite
period family $\mathcal F$ and an actual upper $U_2$ satisfying
$\mathrm{IsUnf}$, with period indices $\mathrm{ik}(k)=k$ and
integral identifications
$\phi_k:\mathcal F.D_k.\mathrm{La}\simeq_{\mathcal O_E}\Lambda_k$.
Every chart satisfies the native $\mathrm{ModelW}$ and the
period, exponential and derivative clauses of $\mathrm{DictLat}$.

\begin{proof}
The coefficient equivalence $e_F$ between the old and upper field
is induced by their equality inside $\mathbb C$. Thus the embedding
of $c_2.F$ given by $(\mathrm{inC}\ c')\circ e_F^{-1}$ is precisely
the inclusion. Its restriction to the level field has the required
root value. Choose a bijection $e$ from $\{1,\ldots,g\}$ to the
pairs $(\tau,a)$ with $a$ ranging over $D_w.\mathrm{mult}(\tau)$
indices. Use the source relabeling for this good embedding and $e$. This gives
the required $\mathrm{IsTypeFn}$ since each relabeled embedding
occurs with its prescribed multiplicity.
The adelic pins, space, uniformization, model and orthogonality laws in $\mathrm{LawsJ}$ give positive
$p+q$, rank $r$, the integral isomorphism $\iota$ onto the identity
lattice in $V$, its alternating trace form, and the given orthogonal
frame $x_0$ at this embedding. These are the fixed
$\mathrm{IotaOK}$ and frame fields used at entry to the cover
construction.

HUF21 applied to this upper moduli scheme gives $\mathcal F$.
Use PUN11 with $G=-V.\mathrm{gram}^{\mathsf T}$,
$\xi_0=-D.\mathrm{host}.\xi_0$, and $\iota$.
In these coordinates the source alternating form is
$\operatorname{Tr}_{E/\mathbb Q}(-\xi_0\mathrm{herm}_G(\iota a,\iota b))$:
interchange the arguments in the lower trace formula and use
alternation, Hermitian symmetry and invariance of trace under
conjugation. The image spans since it is the full identity lattice.
The chosen $\xi_0$ has negative imaginary part on the host CM type.
PUN11 yields on every component the displayed $\phi_k$, the literal
$G,\xi_0$, and positive-sign index $p_0=D_w.q$. Hence the
signature already prescribed for $x_0$ applies on every chart.

For each CM embedding $\tau$ use the homomorphism
$\theta_\tau:\mathbb A_{F,\infty}\otimes_F E\to\mathbb C$
given on elementary tensors by projection to the real place
$\tau|_F$ followed by multiplication by $\tau$ on the second
factor. Its restriction to the first factor is continuous and to the second is $\tau$, giving the original $\mathrm{IsInfEmb}$ restrictions.
Use its restriction to the CM type of each chart as $\mathrm{th}_k$.
Apply NFL17 with D74U2 to obtain a single $\mathrm{FramesOK}$ family
on these models. NFL21 gives $\mathrm{IsUnf}$ for its $U_2$. To
identify the first half of $\mathrm{DictLat}$, the trace identity on
$\phi_k$ is the $\mathrm{ModelBase}$ pairing field. For every
$x\in (U_2.X_k).\mathrm{dom}$ take the period homomorphism with
coordinates $\mathrm{perFr}(U_2.X_k,x,\phi_k a,\mathrm{relab}(e_j)_1,
(e_j)_2)$. The first field of $\mathrm{FramesOK}$ gives its
$\mathrm{IsPeriod}$, the equality of the point with $\mathcal F.u_k(P)$,
and the derivative of $\mathcal F.\mathrm{ex}_k(P)$ on each individual
unit vector indexed by $e^{-1}(\tau,a)$. These are the original
chart identities for all frames in the indicated domain.
\end{proof}
% END D74U3

% BEGIN D74U4
\paragraph{D74U4. A correctly marked chart for every adelic element.}
\label{claim:D74U4}
With the models of D74U3 there is a single dictionary
$h\mapsto(k(h),x(h))$ on $G(\mathbb A_F)$ satisfying the remaining
clauses of $\mathrm{DictLat}$, and every $x(h)$ lies in its indicated
chart domain.

\begin{proof}
Write $O=\mathcal O_E$, $\Lambda=\iota D_w.\mathrm{La}$, and
$N=l_2.M$. For $h\in G(\mathbb A_F)$ take
\[
 \Lambda_h= E^r\cap h_f\widehat\Lambda.
\]
This is a full $O$-lattice with completion $h_f\widehat\Lambda$.
Indeed the two completed lattices are commensurable, agreeing away
from finitely many primes. Choose a positive integer $a$ with
$a\widehat\Lambda\subset h_f\widehat\Lambda\subset
a^{-1}\widehat\Lambda$. The Chinese remainder isomorphism for
$a^{-1}\Lambda/a\Lambda$ identifies its global quotient with the
product of its local quotients. Taking the inverse image of
$h_f\widehat\Lambda/a\widehat\Lambda$ proves the assertion.
The trace pairing of D74U3 is perfect on $\Lambda_h$: at each
rational prime this follows by transport of the perfect pairing
on $\Lambda$ by the unitary $h_f$, and the rational integral dual
is the intersection of its local duals. The alternating pairing,
$O$-action, dimension and signature now define a datum $D_h$ on
$\Lambda_h$. In detail the group is finite free over $\mathbb Z$,
the relation between multiplication by $b$ and by $\bar b$ follows
from the Hermitian pairing, and its real span is the same as that
of $\Lambda$. Hence the original compatible complex structure on
this real vector space has the same positivity and Lie-type traces
on $D_h$.

At every $s>0$ the map $h_f$ induces
$\Lambda/s\Lambda\simeq\Lambda_h/s\Lambda_h$ by completion and
reduction. Choose an $O$-linear lift
$f_s:\Lambda\to\Lambda_h$; it exists since a finite torsion-free
module over the Dedekind domain $O$ is projective. At the unused
index zero set $f_0=0$. For all positive $s,u$ the reductions of
$f_{su}$ and $f_s$ modulo $s$ agree since both reduce $h_f$.
Their pairings agree with the original alternating pairing modulo
$s$ since $h_f$ is unitary, so all pairing multipliers can be one.
The root requirement at auxiliary level $A_2=1$ holds. The combined
modulus is $l_2.M'=\operatorname{lcm}(N,1)=N$, hence the specified
marking is $t_h=f_N\iota$. These data put $(D_h,t_h)$ in the original marked
genus.

By HUF14 and the choice of the family in HUF21, this class has
a representative $\mathcal F.D_k$ and an integral $O$-linear
isometry $u:\Lambda_h\to\mathcal F.D_k.\mathrm{La}$ with
\[
 u f_N(\iota a)-\mathcal F.t_k(a)\ \in N\mathcal F.D_k.\mathrm{La}
 \qquad(a\in\overline{\mathfrak b}_2D_w.\mathrm{La}).
\]
The second marking congruence modulo $A_2=1$ is automatic.
Equivalently one may use PUN15 and PUN22 to construct the point for
this marked class and PUN23 to obtain $u$ from the equality with
its family point. The map $\phi_k u$ extends uniquely to an
$E$-linear automorphism $\gamma$ of $E^r$. It preserves the
Hermitian form: apply equality of trace pairings to $ba,b'$ for
every $b\in O$, then use the perfect field trace pairing and
$\xi_0\ne0$. This puts $\gamma$ in the rational unitary group.
It sends $\Lambda_h$ onto $\Lambda_k$, hence
$\Lambda_k=\Lambda_{\gamma h}$, using rationality of $\gamma$.
The displayed congruence and
$f_N(\iota a)\equiv h_f\iota a\pmod{N\widehat\Lambda_h}$ give
the exact finite part
\[
 \phi_k(\mathcal F.t_k(a))-(\gamma h)_f\iota a\ \in
             N\widehat\Lambda_{\gamma h}
 \qquad(a\in\overline{\mathfrak b}_2D_w.\mathrm{La}).
\]
All rational vectors in this congruence use their finite diagonal
embeddings.

Choose such $k,\gamma$ for each $h$, and set
$x(h)_\tau=\operatorname{archM}_{\theta_\tau}(\gamma h)x_0(\tau)$.
The first factor preserves the chart Gram matrix, so the Gram matrix
of $x(h)_\tau$ is the same signed diagonal as that of $x_0(\tau)$.
Its leading principal minors consequently satisfy all the signed
positivity conditions of the chart, and its determinant is nonzero.
This proves domain membership. The chosen $\gamma$ now supplies
the rational witness, the lattice equality, the finite congruence
and the archimedean equality for this $h$, with the rank
identification $V.\mathrm{rank}=r$ throughout. Together with D74U3
these are all the fields of the precise $\mathrm{DictLat}$.
\end{proof}
% END D74U4

% BEGIN D74U5
\paragraph{D74U5. The original all-tuple cover packet over the prescribed field.}
\label{claim:D74U5}
Fix an original $T,R$, their $\mathrm{LawsJ}$ and level-character
laws, $R.\mathrm{fam\_roots}$, $\mathrm{XiVal}(R)$ and
$\mathrm{CoverField}(T,R)$. For every $i,n$, take any prescribed
$\{I_t,l_{2,t}\}_t$ with the original $\mathrm{IsCoverLevel}$ for
$d_t=\mathrm{tupPrime}(R,t)$ and $\mathrm{HalfIs}$ for the given lower
half ideal. For every coefficient record $c'$ with
$T.Q.c.F\subseteq c'.F=M=\mathrm{embLL}(L)$, these data give
\[
       \mathrm{CoversAt}(T,R,i,n,\mathrm{tupPrime}(R),I,l_2,c',
                              \mathrm{hle}).
\]
On the RTG1 and RTG10 inputs the resulting packet gives the original
$\mathrm{PinsOverL}$ by RTG12.

\begin{proof}
At each $t$, the level fields give $N=md_t$, $(d_t,m)=1$,
$\mathfrak b_2=\mathfrak b I_t$, $I_t+\bar I_t=O$ and $A_2=1$.
The lower auxiliary level divides $A_2$ and is one. From
$\mathfrak b_2\bar{\mathfrak b}_2=(N)$ and
$\mathfrak b\bar{\mathfrak b}=(m)$, cancellation of the nonzero
principal ideal $(m)$ gives $I_t\bar I_t=(d_t)$.
Thus the relative construction A73H1--A73H2 applies with exactly
these same levels. In L73D2 its upper coefficient field can be $M$:
the old coefficient record contains a primitive $m$th root and
$R.\mathrm{fam\_roots}$ gives one for each distinct prime in
$d_t$. Products at coprime orders give a primitive $N$th root.
All its powers belong to $M$, including the specified complex root.
The auxiliary level is one and the original quadratic base is already
in $c'$. The field identification on the actual upper and lower
coefficient records is the identity inside $\mathbb C$. L73D2 now
gives the original $\Delta_t$ at exponent $n$, its level map, framed
lifting, initial character rule and canonical connection.

For each marked prime $P\in R.Q(t)$, the $\mathrm{HalfIs}$ condition
gives $I_t\subseteq P$. The relation $(N)=\mathfrak b_2
\bar{\mathfrak b}_2$ implies $N\in\mathfrak b_2\subseteq I_t$.
Thus the same ramified-place guard as in T73K1 holds for these
prescribed levels as well. Choose its character $\xi_t$ and positive
order bound $m_t$ for the determinants. For any adelic $h$ its
determinant root exists by local Hilbert 90 as in T73K1, so its
value lies in $M$ by $\mathrm{XiVal}$. The arithmetic hypothesis
$\mathrm{CoverField}$ supplies $\mathrm{FnAlgClosed}$ for the
particular constructed upper scheme on this field. On the models and
dictionary of D74U3--D74U4 apply T73K2--T73K3. They produce a
higher rule $f_t$ with $f_t^{m_t}=1$, preserving the full
$\mathrm{IsLev}$ on replacing $\Delta_t.c_\delta$ by $f_t$. In
particular the value formula
\[
 \mathrm{rdU}(U_{2,t}, f_t,\mathrm{dict}_{2,t}(h))
      = \xi_t(\det((\gamma h)_f))
\]
holds for \emph{every} $\gamma$ satisfying
$\mathrm{MarkBy}(h,\gamma)$, including those other than the witness
chosen to define the frame.

Use these tuplewise choices as functions of $t$. The prescribed
integer, ideal and upper level are exactly those of $\Delta_t$.
The level law, $A_2=1$ and the star-ring connection witness give
$\mathrm{CovPin}$. The original half condition and T73K1 give
$\mathrm{HalfIs}$ and $\mathrm{XiIs}$. Take $c_0=1$, of order one,
before testing any tuple in $\mathrm{CXiVal}$. The uniformizations,
component models, index maps, integral isomorphisms, infinite
embeddings and marked formulas just proved give all its fields.
The repeated model clause follows from D74U3, and the last
positive-order clause uses $m_t$. This construction works also when
$d_t=1$ or $n=0$. On the rational-test input RTG12 uses these
same levels, coefficient record and lower frame to fill its eleven
pin fields. The component-field and character-value assumptions in
the displayed cover implication are statements about the original
outer arithmetic data.
\end{proof}
% END D74U5
